\chapter{Судалгааны сэдвийн онол, өнөөгийн түвшин}
% БАГШИД ҮЗҮҮЛЭХ хуулбар (2026-11-10, ажлын хувилбараас). Засварыг Chapters/Chapter2_Literature.tex-д хийнэ.
\label{ch:literature}

%======================================================================
% 2026.11.09-нд хэрэглэгчийн аудитаар өргөтгөсөн. Бүтэц: сургуулийн туслах →
% ресепшн ба олон нийтийн газар → Монгол яриа таних → RAG → төсөөтэй
% системүүдийн харьцуулалт → шинэлэг тал.
% Шинэ ишлэлүүдийн бүтэн текстийг агент уншсан (papers/24-...-20261109).
% BM25, E5, RRF, SQuAD 2.0, Luger-ийн тоо papers/16-reading-notes-45-аас.
% Whisper, Mengke, Lewis-ийг хэрэглэгч өөрөө уншсан (✅).
%======================================================================

\section{Сургуулийн мэдээллийн туслах систем}
\label{sec:lit-campus}

Сүүлийн жилүүдэд их сургуулиуд оюутны асуултад хариулдаг туслах систем хийж байна. Эдгээр нь ихэвчлэн сургуулийн баримтаас хайж, олсон эхэд тулгуурлан хариулдаг. Энэ хэсэгт дөрвөн системийг тоймлов.

\textbf{Unimib Assistant.} Милано-Бикокка их сургуулийн туслахыг ChatGPT-ийн тусгай GPT боломжоор хийсэн \cite{antico2024unimib}. Зохиогчид эхлээд 6 оюутантай ярилцаж хэрэгцээг тодорхойлсон. Оюутнууд дипломын журам, шалгалтын өрөө, анги танхим, хичээлийн хуваарь, тэтгэлгийн мэдээллийг олоход хэцүү гэжээ. Дараа нь өөр 6 оюутнаар туршуулсан. Оюутнууд хурдан, найрсаг хариултыг сайшаасан. Гэхдээ мэдээлэл үргэлж зөв биш, заримдаа байхгүй холбоос зохиосон нь итгэлийг бууруулсан. Оюутнуудын хувьд мэдээллийн эхийг харуулах нь чухал байсан \cite{antico2024unimib}.

\textbf{Marcel.} Марбургийн их сургуулийн туслах элсэгчдийн асуултад хариулдаг \cite{trienes2025marcel}. Баримтын сан нь сургуулийн сайтын 249 баримт ба 36 бэлэн асуулт-хариулт. Асуултыг эхлээд бэлэн асуулттай тулгаж, таарахгүй бол BM25-аар хайдаг. Ийм хайлт BM25 ба векторын хайлтын хослолоос MRR үзүүлэлтээр 75\%-иар илүү гарсан. Хэлний загвар олдсон эх хангалттай бол хариулна, үгүй бол татгалзана. Шалгалтын 95 асуултын 19 нь эх байхгүй асуулт. Зохиогчид олон асуултад систем татгалзсаныг тэмдэглэсэн. Шалтгаан нь сангийн дутуу эсвэл сайтад тэр мэдээлэл огт байхгүй байсан. Үүнийг сайтын дутуу хэсгийг илрүүлэх ашигтай тал гэж үзсэн \cite{trienes2025marcel}.

\textbf{ALOHA.} Бээжингийн их сургуулийн туслах олон хэлээр хариулдаг \cite{tao2025aloha}. Ажлын хуваарь, нээлттэй цаг, зардлын журам зэрэг хүснэгтэн мэдээллийг энгийн текстээс тусад нь хадгалж, асуултын төрлөөр нь аль санг хайхыг шийддэг. Бодит хэрэглэгчийн 100 асуултаар туршихад хятад асуултын 94\%-д, англи асуултын 87\%-д зөв хариулсан. Хайлтгүй ChatGPT ердөө 29\%-д зөв хариулсан. Систем 12 мянга орчим хэрэглэгчид 33 мянга орчим удаа үйлчилсэн \cite{tao2025aloha}.

\textbf{Burch их сургуулийн туслах.} Сайтын 470 гаруй хуудсыг автоматаар татаж, векторын хайлтаар хайдаг \cite{matic2026burch}. 20 асуултаар туршихад бүх асуултын зөв хуудас хайлтын эхний 5-д орсон. Гэхдээ хариултын чанар дунд зэрэг байсан. Асуултын 25\%-д систем зөв хуудсыг олсон ч хариулт агуулсан хэсгийг биш, өөр хэсгийг авсан. Жишээ нь албаны ажлын цагийг асуухад албаны тухай ерөнхий хэсгийг өгсөн \cite{matic2026burch}.

Эдгээр ажлаас гурван зүйл харагдаж байна. Нэгдүгээрт, оюутнууд өрөө, хуваарь зэрэг өдөр тутмын мэдээллийг олоход хүндрэлтэй. Хоёрдугаарт, хайлт сайн байсан ч хариулт дутуу гарах нь элбэг. Гуравдугаарт, эх байхгүй асуултад татгалзах ба эхийг харуулах нь хэрэглэгчийн итгэлд чухал.

\section{Ресепшний туслах ба олон нийтийн газрын ярианы систем}
\label{sec:lit-reception}

Ресепшний туслахын хамгийн сайн баримтжуулсан жишээ бол Карнеги Меллоны их сургуулийн Valerie робот \cite{gockley2005valerie}. Робот барилгын гол хаалганы дэргэдэх лангуунд байрлаж, өрөөний дугаар, зам заах, цаг агаарын мэдээлэл өгдөг. Зохиогчид ресепшний дүрийг сонгосон шалтгаан нь хүмүүс ресепшинтэй яаж харилцахаа аль хэдийн мэддэгт байсан. Есөн сард 16 мянга гаруй удаа хүмүүс роботтой харилцсан, өдөрт дунджаар 88 гаруй хүн. Шинэ зүйлийн сонирхол арилсны дараа нэг харилцаа хагас минут орчим үргэлжилсэн. Энэ нь зам асуухад хангалттай хугацаа юм. Хүмүүсийн бичсэн өгүүлбэрийн 28.1\%-д робот «ойлгосонгүй» гэсэн хариу өгсөн \cite{gockley2005valerie}.

Valerie-д яриа таних биш, гарын товчлуурыг сонгосон. Шалтгаан нь хүн их явдаг хонгилд чимээ, цуурай ихтэй, тухайн үеийн яриа таних систем үүнийг даахгүй байсан. Чихэвч зүүлгэх нь хэрэглэгчид төвөгтэй \cite{gockley2005valerie}. Харин Pappu нар олон хүнтэй амралтын өрөөнд ярианы киоскийг туршсан \cite{pappu2014public}. Систем зөвхөн дуунд тулгуурлахад хүн өөр рүү нь хандаж байгааг 28\%-д л зөв ялгасан. Камераар хүний байрлал, дууны чиглэлийг нэмэхэд зориуд хандсан яриаг 93\%-д зөв ялгасан. Ийм газарт 20 хүний даалгаврын 15-ыг, чимээгүй өрөөнд 16-г амжилттай гүйцэтгэсэн \cite{pappu2014public}. Өөрөөр хэлбэл зөв тохиргоотой бол олон нийтийн газарт ч яриагаар харилцах боломжтой.

Хэрэглэгчийн хүлээлт бас чухал. Luger, Sellen нар ярианы туслах тогтмол хэрэглэдэг 14 хүнтэй ярилцсан \cite{luger2016badpa}. Хэрэглэгчдийн тал нь туслах юу хийж чадахыг мэддэггүй байв. Бүгд нарийн даалгаврын үр дүнг дэлгэцээр харж шалгахыг хүссэн. Иймээс дуугаар хариулахдаа хариултыг дэлгэцэнд эхтэй нь харуулах хэрэгтэй.

\section{Монгол хэлний яриа таних}
\label{sec:lit-asr}

Яриа таних (\ac{ASR}) гэдэг нь хүний ярианы аудиог текст болгох ажил юм. Манай системд хүн асуултаа ярихад эхлээд энэ алхам ажиллана. Энд алдаа гарвал систем буруу мэдээлэл хайна. Яриа танилтын чанарыг ихэвчлэн үгийн алдааны хувиар (\ac{WER}) хэмждэг:
\begin{equation}
\mathrm{WER} = \frac{S + D + I}{N} \times 100\%
\label{eq:wer}
\end{equation}
Энд $S$ нь буруу таньсан үг, $D$ нь орхисон үг, $I$ нь илүү нэмсэн үгийн тоо, $N$ нь зөв бичвэрийн нийт үгийн тоо. Тэмдэгтийн алдааны хувийг (\ac{CER}) мөн ийм аргаар үсгээр тооцдог.

\subsection{Whisper}

Whisper бол OpenAI-ийн нээлттэй яриа таних загвар \cite{radford2023whisper}. Үүнийг интернэтээс авсан 680 мянган цагийн аудио ба бичвэрээр сургасан. Үүний 117 мянган цаг нь англиас бусад 96 хэлнийх. Whisper-ийг нэмэлт сургалтгүйгээр шууд хэрэглэж болдог. Гэхдээ сургалтын өгөгдөлд Монгол яриа цөөн байсан. Иймээс Whisper large-v2 загварын Монгол хэлний WER FLEURS өгөгдөл дээр 110.5\%, Common Voice 9 дээр 128.4\% гарсан \cite{radford2023whisper}. % Хавсралтын Table 13 (FLEURS), Table 11 (CV9)-тэй агент тулгасан (11/06).
Мөн урт аудиод үг давтах, огт хэлээгүй үг бичих алдаа гардаг \cite{radford2023whisper}.

Энэ ажилд Монгол хэлээр нэмэлт сургасан Whisper small загварыг (\texttt{Cafet/\allowbreak whisper-small-\allowbreak mongolian}) харьцуулалтын суурь болгосон \cite{cafet_whisper_small_mn}. Энэ загвар өөрийн компьютер дээр ажиллана.

\subsection{Whisper-ийг Монгол хэлэнд сайжруулах судалгаа}

Mengke нар (2025) Халх Монгол хэлний яриа танилтыг сайжруулах арга санал болгосон \cite{mengke2025khalkha}. Common Voice 17-д Халх Монгол хэлний сургалтын өгөгдөл ердөө 3 цаг байдаг. Иймээс зохиогчид Өвөр Монголын Чахар аялгууны 147 цагийн өгөгдлийг кирилл бичигт хөрвүүлсэн. Whisper large-v3 загварыг шууд хэрэглэхэд WER 92.60\% байсан. Халхын 3 цагийн өгөгдлөөр сургахад 43.02\% болсон. Эхлээд Чахарын өгөгдлөөр, дараа нь Халхын өгөгдлөөр сургахад 31.45\% хүртэл буурсан \cite{mengke2025khalkha}. Өөрөөр хэлбэл алдааны ихэнх бууралт Халхын цөөн өгөгдлөөр сургахад гарсан. Энэ судалгаа Whisper-ийг Монгол хэлэнд сайжруулж болохыг харуулж байна. Гэхдээ үүнд өгөгдөл, хүчтэй \ac{GPU} хэрэгтэй.
% Тоог өгүүллийн 4-р хэсгийн хүснэгттэй агент тулгасан (11/06).

Том хэлний загварууд ч Монгол хэлэнд сул хэвээр. MM-Eval нь орчин үеийн кирилл Монгол хэлээр 1840 даалгавартай шалгалт \cite{zhang2024mmeval}. Үүгээр туршихад таван загвар бүгд өгүүлбэрийн дүрмийн даалгаварт утгын даалгавраас илүү гарсан. Хамгийн сайн GPT-4 Turbo дүрэмд 90.69\%, утгад 72.53\% авсан. Зохиогчид загварууд Монгол хэлний гүн утгыг ойлгоход сул гэж дүгнэсэн \cite{zhang2024mmeval}. Энэ нь манай системд хариултыг загвараар шинээр бичүүлэхгүй, эхийн үгээр өгөх шийдвэрийг дэмжинэ.

\subsection{Gemini 3.5 Transcribe}

Gemini 3.5 Transcribe бол Google-ийн яриаг текст болгох онлайн үйлчилгээ \cite{google_transcribe_model}. Монгол хэлийг (\texttt{mn-MN}) дэмждэг. Үүнийг \ac{API} түлхүүрээр ашиглана. Аудио Google-ийн сервер рүү очдог тул интернэт хэрэгтэй. Үнэгүй квоттой, төлбөртэй хувилбар нь минутад 0.005 ам.доллар орчим \cite{google_pricing}. Энэ үйлчилгээний Монгол хэлний алдааг хэмжсэн судалгаа олдоогүй. Иймээс бид өөрсдийн бичлэг дээр хэмжсэн (§\ref{sec:res-stt}). \tref{tab:asr_compare}-т хоёр аргын гол ялгааг харуулав.

\begin{table}[ht]
\centering
\caption{Whisper ба Gemini 3.5 Transcribe-ийн ялгаа}
\label{tab:asr_compare}
\begin{tabular}{|p{3.6cm}|p{5.0cm}|p{5.0cm}|}
\hline
 & \textbf{Whisper (Cafet small)} & \textbf{Gemini 3.5 Transcribe} \\
\hline
Хаана ажиллах & Өөрийн компьютер & Google-ийн сервер \\
\hline
Интернэт & Хэрэггүй & Хэрэгтэй \\
\hline
Аудио & Компьютер дотроо үлдэнэ & Google руу очно \\
\hline
Сайжруулах & Нэмэлт сургалт хэрэгтэй & Загварыг сургах боломжгүй, нэр томьёоны жагсаалт өгч болно \\
\hline
Зардал & Компьютерын нөөц & Үнэгүй квот, дараа нь төлбөртэй \\
\hline
\end{tabular}
\end{table}

\subsection{Яриа танилтын алдаа хайлтад хэрхэн нөлөөлөх вэ}

Дуугаар асуухад яриа танилтын алдаа дараагийн алхам болох хайлт руу дамждаг. Sidiropoulos нар англи асуултыг дуу болгоод буцааж текст болгосноор ийм алдааг дуурайлгасан \cite{sidiropoulos2022asr}. Алдаатай асуултад BM25-ын хайлтын үзүүлэлт (MRR@10) 62.66\%-иар, векторын хайлтынх 57.72\%-иар буурсан. Жинхэнэ хүмүүсийн, янз бүрийн аялгатай яриа дуурайлгасан алдаанаас ч хэцүү байсан. Хамгийн элбэг алдаа нь үгийг буруу хуваах, нэрийг дуудлагаар нь буруу бичих байсан \cite{sidiropoulos2022asr}. Сургуулийн асуулт нэр, тоо ихтэй тул энэ нь манай системд шууд хамаатай.

\section{Хайлтаар баяжуулсан хариулт (RAG)}
\label{sec:lit-rag}
% Lewis нар 2020 (RAG01) — хэрэглэгч 11/07-нд уншсан.

Хэлний том загвар сургалтын үедээ олон баримтыг жин дотроо хадгалдаг. Lewis нар үүнийг \textit{жиндээ хадгалсан мэдлэг} гэж нэрлэсэн \cite{lewis2020rag}. Ийм мэдлэгийг шинэчлэхэд загварыг дахин сургах хэрэгтэй. Загвар хариултаа аль эхээс авсныг харуулж чаддаггүй, заримдаа байхгүй баримт зохиодог \cite{lewis2020rag}.

RAG нь загварт \textit{гадаад санах ой} нэмдэг. Энэ нь тусдаа хадгалсан бичвэрийн сан юм. Систем хоёр хэсэгтэй \cite{lewis2020rag}:
\begin{enumerate}
    \item \textbf{Хайгч} асуултад хамгийн ойр $K$ баримтыг сангаас олно. Lewis нар DPR хайгчийг ашигласан. Асуулт ба баримтыг тоон вектор болгоод, векторын ойролцоогоор эрэмбэлнэ.
    \item \textbf{Бичигч} асуулт ба олдсон баримтыг хараад хариулт бичнэ. Үүнд 400 сая орчим параметртэй BART загварыг ашигласан.
\end{enumerate}
Баримтын сан нь 2018 оны 12-р сарын Wikipedia байсан. Нийтлэлүүдийг 100 үгтэй хэсгүүдэд хувааж, 21 сая баримт болгосон \cite{lewis2020rag}.

Нээлттэй асуулт-хариултын Natural Questions өгөгдөл дээр RAG 44.5\% яг таарсан хариулт авсан. Зөвхөн жиндээ мэдлэгтэй T5-11B+SSM загвар 36.6\% авсан \cite{lewis2020rag}. Мөн зохиогчид векторын хайгчийг BM25-аар сольж туршсан. Ихэнх ажилд векторын хайгч илүү, харин хүн, газрын нэр их агуулсан FEVER өгөгдөл дээр BM25 илүү байсан \cite{lewis2020rag}. 2016 ба 2018 оны Wikipedia-гаар баримтын санг солиход 2018 оны улсын удирдагчийн асуултын зөв хариулт 4\%-иас 68\% болсон. Өөрөөр хэлбэл загварыг дахин сургахгүйгээр зөвхөн санг солиод мэдлэгийг шинэчилж болно \cite{lewis2020rag}.

\subsection{Хариултыг бичих үү, эхээс сонгох уу}
\label{sec:lit-selective}

Манай систем RAG-ийн хайгч хэсгийг ашигладаг ч бичигч загваргүй. Хариултыг шинээр бичихгүй, олдсон хэсгүүдээс дүрмээр нэгийг сонгоод эхийн үгээр нь өгнө. Иймээс манай систем хариултыг эхийн текстээс сонгож авдаг асуулт-хариултын аргад илүү ойр. SQuAD өгөгдөлд ч хариулт нь өгөгдсөн догол мөрийн нэг хэсэг байдаг \cite{rajpurkar2018squad2}. Gaglio нарын кампусын туслах ч баримтын сангаас хамгийн ойр хэсгүүдийг олж, эх баримтын нэртэй нь буцаадаг \cite{gaglio2019campus}.

Эхийн үгээр хариулахын давуу тал нь хариултад эхэд байхгүй үг нэмэгдэхгүй. Гэхдээ эх заасан хариулт зөв хариулт гэсэн үг биш. Gao нар хариултыг эх заан бичдэг системийг хэмжих ALCE аргыг гаргасан \cite{gao2023alce}. Хамгийн дээр олдсон хэсгийг хариулт болгон шууд хуулахад эх заалтын үзүүлэлт 99.4 болсон ч асуултын зөв хариултын үзүүлэлт 35.1 л байсан \cite{gao2023alce}. Өөрөөр хэлбэл эх нь зөв заагдсан ч тэр эх асуултад хариулж байгаа эсэхийг тусад нь шалгах хэрэгтэй. Иймээс бид хариулт бүрийг «шаардсан утга гарсан эсэх»-ээр тусад нь хэмжсэн (\ref{sec:res-answers}).

\subsection{Үгээр ба утгаар хайх}

BM25 нь асуултын үг баримтад хэр олон удаа, хэр ховор үг байгаагаар оноо өгдөг \cite{robertson2009bm25}. Үг олон давтагдах тусам оноо өсөх боловч тодорхой хэмжээнээс хэтэрдэггүй. Урт баримтын оноог багасгаж засдаг. BM25 үгийн дарааллыг тооцдоггүй, үгсийг хоорондоо хамааралгүй гэж үздэг \cite{robertson2009bm25}.

Утгаар хайхад текстийг тоон вектор болгодог загвар хэрэглэнэ. Multilingual E5 загварыг олон хэлний тэрбум орчим текстийн хосоор сургасан \cite{wang2024me5}. 18 хэлний MIRACL өгөгдөл дээр жижиг E5 загвар nDCG@10 үзүүлэлтээр 60.8, BM25 39.3 авсан. Гэхдээ энэ үнэлгээнд Монгол хэл ороогүй \cite{wang2024me5}.

Хоёр хайлтын жагсаалтыг нэгтгэх энгийн арга бол \ac{RRF} \cite{cormack2009rrf}. Баримт бүрт жагсаалт бүрээс $1/(k + \text{байр})$ оноо нэмнэ. Зохиогчид $k=60$ ашигласан. Энэ арга оноог биш, зөвхөн байрыг ашигладаг тул өөр хэмжүүртэй хайлтуудыг шууд нэгтгэж болно. TREC өгөгдөл дээр RRF хамгийн сайн дан системээс илүү гарсан: MAP .2145 ба .2016 \cite{cormack2009rrf}.

\subsection{Хариулт байхгүй асуулт}

Хариулт нь эхэд байхгүй асуултад систем таамаглаж хариулах нь аюултай. Rajpurkar нар ийм 53 мянга гаруй асуулттай SQuAD 2.0 өгөгдлийг гаргасан \cite{rajpurkar2018squad2}. Асуултууд эхтэй холбоотой, үнэмшилтэй харагддаг боловч хариулт нь эхэд байхгүй. Тухайн үеийн хамгийн сайн загвар F1 66.3 авсан бол хүн 89.5 авсан \cite{rajpurkar2018squad2}. Өөрөөр хэлбэл «мэдэхгүй» гэж зөв хэлэх нь хариулт олохоос ч хэцүү. Marcel системд ч эх байхгүй асуултыг тусад нь хэмжсэн \cite{trienes2025marcel}. Иймээс манай шалгалтын багцад эх байхгүй асуултыг тусад нь оруулсан.

\subsection{Хариулж чадахгүй үед юу хийх вэ}
\label{sec:lit-referral}

Хариулт олдохгүй үед систем «мэдэхгүй» гээд зогсох нь хэрэглэгчид тус болохгүй. Ashktorab нар 203 оролцогчтой туршилтаар чатбот ойлгоогүй үед ойролцоо 2--3 сонголт санал болгохыг хэрэглэгчид хамгийн их илүүд үзсэнийг харуулсан \cite{ashktorab2019resilient}. Хүн ажилтан руу шилжүүлэхийг хэрэглэгч өөрөө дахин оролдоод бүтэлгүйтсэн үед л илүүд үзсэн \cite{ashktorab2019resilient}. Følstad, Taylor нар бодит үйлчилгээний чатботын яриаг шинжилж, «итгэлгүй байна» гэж хэлээд өөр сонголт санал болгодог болгосны дараа буруу хариулт 28\%-иас 14\% болж буурсныг тайлагнасан \cite{folstad2021investigating}. Хотын үйлчилгээний чатботын хэрэглэгчид ерөнхий холбоосоос илүү тодорхой албыг заахыг хүссэн \cite{zhou2025publicservice}. Эдгээр нь өөр хэл, өөр орчны судалгаа боловч манай систем хариулт олдохгүй үед хаанаас лавлахыг заах шийдвэрийг (\ref{sec:requirements}) дэмжиж байна.

\subsection{Мэдээллийн шинэлэг байдал}

Сургуулийн мэдээлэл байнга өөрчлөгддөг. Ouyang нар хуучирсан мэдээлэл RAG системд нөлөөлөхийг шалгаж, хайлт хуучин ба шинэ хувилбарыг бараг адил сайн олдгийг, хуучирсан мэдээлэл огт мэдээлэлгүй байхаас ч илүү хор хөнөөлтэй болохыг харуулсан \cite{ouyang2025hoh}. Иймээс сайтын хуулбарыг тогтмол шинэчилж, хуучин хувилбарыг хайлтаас хасах, хариулт бүрт огноо харуулах шаардлагатай.

\section{Төсөөтэй системүүдийн харьцуулалт}
\label{sec:lit-compare}

\tref{tab:similar_systems}-т дээрх системүүдийг манай системтэй харьцуулав.

\begin{table}[H]
\centering
\caption{Төсөөтэй системүүдийн харьцуулалт}
\label{tab:similar_systems}
\small
\begin{tabular}{|p{2.2cm}|p{1.7cm}|p{1.6cm}|p{3.2cm}|p{2.4cm}|p{3.0cm}|}
\hline
\textbf{Систем} & \textbf{Хэл / газар} & \textbf{Асуух} & \textbf{Мэдээллийн эх} & \textbf{Хариулт} & \textbf{Үнэлгээ} \\
\hline
Unimib \cite{antico2024unimib} & Итали их сургууль & Бичгээр & Баримт, холбоос & ChatGPT бичнэ & 6 оюутны туршилт \\
\hline
Marcel \cite{trienes2025marcel} & Герман их сургууль & Бичгээр & Сайтын 249 баримт, 36 бэлэн хариулт & Загвар бичнэ, эхгүй бол татгалзана & 95 асуулт (19 эхгүй) \\
\hline
ALOHA \cite{tao2025aloha} & Хятад, англи, франц & Бичгээр & Гарын авлага, сайт, хүснэгт & GPT-4o бичнэ & Бодит 100 асуулт \\
\hline
Burch \cite{matic2026burch} & Англи & Бичгээр & Сайтын 470+ хуудас & Загвар бичнэ & 20 асуулт \\
\hline
Valerie \cite{gockley2005valerie} & Англи & Гараар & Өрөө, зам, цаг агаар & Бэлэн хариулт & 9 сарын ашиглалт \\
\hline
Gaglio нар \cite{gaglio2019campus} & Итали их сургууль & Вэб чат, апп & 50 орчим баримт (журам, FAQ, сайт) & Бэлэн хариулт эсвэл эх нэртэй хэсэг & Хэрэглэгчийн бүлэг \\
\hline
\textbf{Манай} & \textbf{Монгол} & \textbf{Дуугаар, бичгээр} & \textbf{Сайт (75 хуудас); сургалтын албаны хүснэгт бэлтгэгдэж байна} & \textbf{Эхийн үгээр, эхгүй бол татгалзана} & \textbf{Тусдаа: 12 бичлэг; бичгийн 40 асуулт} \\
\hline
\end{tabular}
\end{table}

Хүснэгтийн ихэнх систем хариултыг хэлний загвараар шинээр бичүүлдэг. Манай систем эхийн үгийг өөрчлөхгүй өгдөг. Ингэснээр хариултад эхэд байхгүй баримт нэмэгдэхгүй, хариулт бүрийг эхтэй нь шууд тулгаж болно. Монгол хэлээр хариулт бичих загварын чанарыг энэ ажилд хэмжээгүй. Хүснэгтийн сургуулийн туслахууд бүгд бичгээр асуудаг. Gaglio нар яриаг текст болгох модулийг ашиглаж болохыг дурдсан ч хэрэгжүүлсэн интерфейс нь вэб чат ба гар утасны апп байсан \cite{gaglio2019campus}. Манай систем ресепшнд байрлах тул дуугаар асуух боломжтой. Гэхдээ дуугаар асуух бүтэн урсгалыг энэ ноорогт хэмжээгүй (5.3).

\section{Судалгааны ажлын шинэлэг тал}
\label{sec:lit-novelty}

Бид 2026 оны 11-р сард вэб хайлтаар өмнөх ажлыг хайсан. Түлхүүр үг нь «university campus chatbot», «robot receptionist», «voice assistant kiosk public space», «Mongolian question answering», «RAG university chatbot» байв. Энэ нь системтэй тойм биш. Иймээс доорх зүйлсийг «огт байхгүй» гэж биш, бидний хайлтад олдоогүй хослол гэж үзнэ. Тухайлбал Gaglio нарын ажил энэ хайлтаар олдоогүй, дараа нь нэмэгдсэн.

\begin{enumerate}
    \item \textbf{Монгол хэл ба ресепшн.} Сургуулийн туслах систем олон бий (\ref{sec:lit-compare}). Манай ялгаа нь Монгол хэлээр, сургуулийн ресепшнд зориулж, дуугаар болон бичгээр асуух боломжтой хийсэн явдал.
    \item \textbf{Ижил бичлэг дээрх харьцуулалт.} Онлайн яриа таних үйлчилгээ (Gemini 3.5 Transcribe) ба Монгол хэлээр сургасан Whisper загварыг сургуулийн асуултын ижил бичлэг дээр харьцуулсан. Энэ нь нэг хүний бичлэгээр хийсэн жижиг хэмжилт.
    \item \textbf{Монгол сургуулийн мэдээлэлд хэрэглэсэн.} Эхээс хэсэг буцаах \cite{gaglio2019campus}, эх байхгүй асуултыг тусад нь хэмжих \cite{rajpurkar2018squad2,trienes2025marcel} аргууд өмнө нь бий. Бид эдгээрийг Монгол сургуулийн мэдээлэлд хэрэглэж, урьдчилж түгжсэн багцаар хэмжсэн.
    \item \textbf{Мэдлэгийн мод.} Системийн мэдэхгүй байгаа хэсгийг сэдвээр нь илрүүлж, сургуулиас дутуу мэдээллийг хүсэхэд ашигласан.
\end{enumerate}
